% ----------------------------------------------------------
\chapter{Conclusão}\label{cap:conclusao}
% ----------------------------------------------------------

A aplicação do método de Monte Carlo evidenciou de forma clara a intersecção entre processos probabilísticos e o cálculo de áreas geométricas determinísticas. Através do desenvolvimento e da execução do algoritmo, confirmou-se que a constância da aproximação obedece à Lei dos Grandes Números. 

Observou-se que amostras computacionais limitadas sujeitam os resultados a flutuações e ruídos numéricos transientes. Contudo, ao expandir a malha de sorteios para escalas na ordem de milhões, o erro relativo percentual diminuiu de forma drástica, garantindo convergência com o valor tabelado de $\pi \approx 3,14159$. 

Este estudo de caso reforça não apenas a aplicabilidade da técnica na teoria de integração numérica multidimensional, mas também a necessidade de compreender as limitações dos geradores de números pseudoaleatórios na confiabilidade científica dos dados.